\documentclass[11pt]{amsart}
\usepackage{amsmath,amssymb,amsthm}
\usepackage[margin=1in]{geometry}
\usepackage{hyperref}

\newtheorem{theorem}{Theorem}[section]
\newtheorem{corollary}[theorem]{Corollary}
\newtheorem{lemma}[theorem]{Lemma}
\newtheorem{conjecture}[theorem]{Conjecture}
\theoremstyle{remark}
\newtheorem{remark}[theorem]{Remark}
\newtheorem{example}[theorem]{Example}

\newcommand{\Ga}{\mathcal{G}}
\newcommand{\cF}{\mathcal{F}}
\newcommand{\cA}{\mathcal{A}}
\newcommand{\cH}{\mathcal{H}}
\newcommand{\N}{\mathbb{N}}
\newcommand{\R}{\mathbb{R}}
\newcommand{\Q}{\mathbb{Q}}
\newcommand{\Pb}{\mathbb{P}}
\newcommand{\E}{\mathbb{E}}
\newcommand{\esssup}{\operatorname*{ess\,sup}}
\newcommand{\dist}{\operatorname{dist}}

\title[Non-endogeny for discounted tree sums]{Non-endogenous fixed points of the discounted tree sum equation: the regular case}
\author{Yitzchak Shmalo}
\date{}

\begin{document}

\begin{abstract}
Aldous and Bandyopadhyay showed that every nonzero solution $Y$ of the homogeneous
equation $Y\stackrel{d}{=}\max_i \xi_i Y_i$ produces, for each $a\ge 0$, an invariant
measure $\mu_a$ of the discounted tree sum equation
$X\stackrel{d}{=}\eta+\max_i \xi_i X_i$, and conjectured that for $a>0$ the invariant
recursive tree process with marginal $\mu_a$ is never endogenous. We show that the
conditional essential supremum of $Y$ given all the innovations is an endogenous solution
of the homogeneous equation with values in $[0,\infty]$, and that $\mu_a$ fails to be
endogenous as soon as this supremum is infinite with positive probability. It follows
that the conjecture, for all nonzero $Y$ at once, is equivalent to the absence of a
nonnull finite endogenous solution of the homogeneous equation. When the weights have a
characteristic exponent, the disintegration theorem of Alsmeyer, Biggins and Meiners
shows that the conditional supremum is infinite on the survival event of an endogenous
fixed point of the smoothing transform, and the conjecture holds.
\end{abstract}

\maketitle

\section{Introduction}

Let $(\eta,(\xi_i)_{i\ge 1})$ be a random element of $[0,\infty)\times[0,\infty)^{\N}$ and
consider the recursive distributional equation
\begin{equation}\label{eq:rde}
X\stackrel{d}{=}\eta+\sup_{i\ge 1}\xi_i X_i \qquad\text{on } [0,\infty),
\end{equation}
where $X_1,X_2,\dots$ are independent copies of $X$, independent of $(\eta,(\xi_i))$.
Equation \eqref{eq:rde} is the discounted tree sum equation of Aldous and Bandyopadhyay
\cite[Section~4.4]{AB}. On the tree of finite words over $\N$, place independent copies
$(\eta_v,(\xi_{vi})_{i\ge1})$ of the innovation at the vertices $v$, and let $L(v)$ be the
product of the weights along the path from the root $\varnothing$ to $v$, with
$L(\varnothing)=1$. The variable
\begin{equation}\label{eq:X}
X=\sup_{\varnothing=v_0,v_1,v_2,\dots}\ \sum_{d\ge 0}\eta_{v_d}L(v_d),
\end{equation}
the supremum being over rays from the root, is a solution of \eqref{eq:rde} whenever it is
finite almost surely, and its law $\mu$ is the lower invariant measure. Dropping $\eta$
gives the homogeneous equation
\begin{equation}\label{eq:hom}
Y\stackrel{d}{=}\sup_{i\ge 1}\xi_i Y_i\qquad\text{on } [0,\infty).
\end{equation}
Aldous and Bandyopadhyay proved the following \cite[Corollary~33]{AB}: if $X$ in
\eqref{eq:X} is finite almost surely and $\dist(Y)$ is a nonzero solution of
\eqref{eq:hom}, then for every $a\ge 0$ the iterates of the map of \eqref{eq:rde} started
at $\dist(aY)$ converge to an invariant measure $\mu_a$; if moreover
\begin{equation}\label{eq:39}
\eta \text{ is independent of } (\xi_i), \text{ and } 0 \text{ is in the support of } \eta,
\end{equation}
then the measures $\mu_a$ are distinct. They stated the following.

\begin{conjecture}[{\cite[Conjecture~34]{AB}}]\label{conj}
Under the assumptions of Corollary~33 of \cite{AB} and \eqref{eq:39}, for $a>0$ the
invariant recursive tree process associated with $\mu_a$ is not endogenous.
\end{conjecture}

Endogeny means that the value at the root is a measurable function of the innovations
\cite[Definition~7]{AB}; the precise setting is recalled in Section~\ref{sec:prelim}.
Let $\Ga$ be the $\sigma$-field generated by all the innovations
$(\eta_v,(\xi_{vi})_i)$, and let $(Y_v)$ be the invariant recursive tree process of
\eqref{eq:hom} with marginal $\dist(Y)$, built on the same tree. Put
\[
M=\esssup\,(Y_\varnothing\mid \Ga)\in[0,\infty].
\]
Our first result reduces the conjecture to a statement about \eqref{eq:hom} alone.

\begin{theorem}\label{thm:main}
Suppose $X$ in \eqref{eq:X} is finite almost surely and let $\dist(Y)$ be a nonzero
solution of \eqref{eq:hom}.
\begin{enumerate}
\item[(a)] $M$ is an endogenous solution of \eqref{eq:hom} on $[0,\infty]$, and
$M\ge Y_\varnothing$ almost surely.
\item[(b)] If the recursive tree process of $Y$ is endogenous, then $\mu_a$ is endogenous
for every $a\ge0$.
\item[(c)] If $\Pb(M=\infty)>0$, then $\mu_a$ is not endogenous for any $a>0$.
\end{enumerate}
Consequently the following are equivalent:
\begin{enumerate}
\item[(i)] for every nonzero solution $\dist(Y)$ of \eqref{eq:hom} and every $a>0$, the
measure $\mu_a$ is not endogenous;
\item[(ii)] the only endogenous solution of \eqref{eq:hom} on $[0,\infty)$ is $\delta_0$.
\end{enumerate}
\end{theorem}

Condition \eqref{eq:39} plays no role in Theorem~\ref{thm:main}. Our second result settles
Conjecture~\ref{conj} when the weights have a characteristic exponent. Let
$N=\#\{i:\xi_i>0\}$ and
\[
m(\theta)=\E\sum_{i:\,\xi_i>0}\xi_i^{\theta}\in[0,\infty],\qquad \theta\ge0 .
\]

\begin{theorem}\label{thm:regular}
Suppose $X$ in \eqref{eq:X} is finite almost surely, that $1<m(\theta_0)<\infty$ for some
$\theta_0\ge 0$ and $m(p)<1$ for some $p>\theta_0$, and that either $N<\infty$ almost surely
or $\E\, W_1\log^+W_1<\infty$, where $W_1=\sum_i\xi_i^{\alpha}$ and $\alpha$ is the unique
point of $(\theta_0,p)$ with $m(\alpha)=1$. Then for every nonzero solution $\dist(Y)$ of
\eqref{eq:hom} we have $\Pb(M=\infty)>0$, and for every $a>0$ the invariant recursive tree
process associated with $\mu_a$ is not endogenous.
\end{theorem}

The proof identifies $M$ through the disintegration of fixed points of the smoothing
transform due to Alsmeyer, Biggins and Meiners \cite[Theorem~8.3]{ABM}: conditionally on
the weights, $1/Y_\varnothing$ has a Weibull-type survival function
$\exp(-Wh(t)t^{\alpha})$ with $W$ an endogenous fixed point of the smoothing transform with
weights $(\xi_i^{\alpha})$, so $Y_\varnothing$ is conditionally unbounded on $\{W>0\}$.
Theorem~\ref{thm:regular} covers Examples~37 and~38 of \cite{AB}; see
Section~\ref{sec:remarks}, which also discusses what remains open.

\section{Preliminaries}\label{sec:prelim}

\subsection{Recursive tree processes}
Let $V=\bigcup_{n\ge0}\N^n$ be the set of finite words, with root $\varnothing$, children
$vi$ of $v$, and $|v|$ the length of $v$. Let $\theta_v=(\eta_v,(\xi_{vi})_{i\ge1})$,
$v\in V$, be independent copies of $(\eta,(\xi_i))$. For $v\in V$ let $\Ga_v$ be the
$\sigma$-field generated by $(\theta_w: w\in vV)$, the innovations of the subtree rooted at
$v$, so that $\Ga=\Ga_\varnothing$. For $w\in vV$ write $L_v(w)$ for the product of the
weights along the path from $v$ to $w$, with $L_v(v)=1$, and $L=L_\varnothing$.

Consider a recursive distributional equation of the form $Z\stackrel{d}{=}g(\theta;Z_1,Z_2,
\dots)$. An invariant recursive tree process with marginal $\nu$ is a family
$(Z_v)_{v\in V}$, defined jointly with $(\theta_v)$, such that
$Z_v=g(\theta_v;Z_{v1},Z_{v2},\dots)$ almost surely for every $v$, and such that for every
$n$ the variables $(Z_v:|v|=n)$ are independent with law $\nu$ and independent of
$(\theta_w:|w|<n)$. By \cite[Lemma~6]{AB} such a process exists if and only if $\nu$ is a
solution. Its joint law with the innovations is determined by $\nu$: it is determined by
the laws of the families $(\theta_w:|w|<k;\ Z_w:|w|=k)$, $k\ge0$, which are product laws,
since every $Z_w$ with $|w|<k$ is almost surely a function of such a family. The process
is endogenous if $Z_\varnothing$ is $\Ga$-measurable (up to a null set); by the uniqueness
just noted this is a property of $\nu$, and we say that $\nu$ is endogenous.

\begin{lemma}\label{lem:blocks}
Let $(Z_v)$ be an invariant recursive tree process and, for $v\in V$, let
$\cF_v$ be the $\sigma$-field generated by $(\theta_w,Z_w:w\in vV)$. For every $n$, the
$\sigma$-fields $\cF_v$, $|v|=n$, and $\sigma(\theta_w:|w|<n)$ are independent, and the
family $(\theta_{vw},Z_{vw})_{w\in V}$ has the same law as $(\theta_w,Z_w)_{w\in V}$ for
every $v$.
\end{lemma}

\begin{proof}
For $k\ge n$, the family $(\theta_w:|w|<k;\ Z_w:|w|=k)$ consists of independent variables.
Grouping it according to the vertex of depth $n$ above each $w$ with $|w|\ge n$ gives
independent blocks, independent of $(\theta_w:|w|<n)$, and the block of $v$ has the law of
the family $(\theta_w:|w|<k-n;\ Z_w:|w|=k-n)$. The variables $Z_w$ with $n\le|w|\le k$ are
almost surely functions of these blocks. As $k\to\infty$ the $\sigma$-fields generated by
the blocks increase and generate $\cF_v$ up to null sets, and the claims follow.
\end{proof}

We use the process of \eqref{eq:hom} with innovations $\theta_v$; the map of \eqref{eq:hom}
does not use $\eta_v$.

\begin{lemma}\label{lem:eta}
Let $(Y_v)$ be the invariant recursive tree process of \eqref{eq:hom} with innovations
$(\theta_v)$. Conditionally on $\sigma(\xi_{vi}:v\in V,i\ge1)$, the families $(Y_v)$ and
$(\eta_v)$ are independent. In particular
$\Pb(Y_\varnothing>y\mid\Ga)=\Pb(Y_\varnothing>y\mid \xi_{vi}: v\in V, i\ge 1)$ almost
surely for every $y$.
\end{lemma}

\begin{proof}
Take an invariant recursive tree process $(Y_v)$ of \eqref{eq:hom} with innovations
$(\xi_{vi})_i$ only, and draw each $\eta_v$ from the conditional law of $\eta$ given
$(\xi_i)=(\xi_{vi})_i$, independently over $v$ and of $(Y_v)$ given the weights. The
pairs $\theta_v$ are then independent copies of $(\eta,(\xi_i))$, and for every $n$ the
variables $(Y_v:|v|=n)$ are independent of $(\theta_w:|w|<n)$, since $(\eta_w:|w|<n)$ is a
function of $(\xi_{wi}:|w|<n)$ and of randomness independent of everything else. So this
is an invariant recursive tree process with innovations $(\theta_v)$, it has the stated
conditional independence, and by uniqueness in law every such process has it.
\end{proof}

\subsection{Conditional essential supremum}
For a $[0,\infty]$-valued random variable $U$ and a $\sigma$-field $\cH$, the conditional
essential supremum $\esssup(U\mid\cH)$ is the $\cH$-measurable $[0,\infty]$-valued
variable $S$, unique up to null sets, such that $U\le S$ almost surely and $S\le S'$
almost surely for every $\cH$-measurable $S'$ with $U\le S'$ almost surely. One may take
$S=\sup\{q\in\Q_+:\Pb(U>q\mid\cH)>0\}$ with $\sup\varnothing=0$. We use three elementary
facts, valid almost surely:
\begin{enumerate}
\item[(E1)] $\esssup(\sup_i U_i\mid\cH)=\sup_i\esssup(U_i\mid\cH)$ for countably many
$U_i\ge0$;
\item[(E2)] $\esssup(cU\mid\cH)=c\,\esssup(U\mid\cH)$ for $\cH$-measurable
$c\in[0,\infty)$, with $0\cdot\infty=0$;
\item[(E3)] if $\sigma(U)\vee\cH_1$ is independent of $\cH_2$, then
$\esssup(U\mid\cH_1\vee\cH_2)=\esssup(U\mid\cH_1)$.
\end{enumerate}
(E1) and (E2) follow from the definition, and (E3) from
$\Pb(U>q\mid\cH_1\vee\cH_2)=\Pb(U>q\mid\cH_1)$.

\section{The conditional supremum and the coupling}\label{sec:proof}

Throughout this section $X$ in \eqref{eq:X} is finite almost surely, $\dist(Y)$ is a
nonzero solution of \eqref{eq:hom}, and $(Y_v)$ is its invariant recursive tree process
with innovations $(\theta_v)$. Fix a measurable $m_0$ with
$\esssup(Y_\varnothing\mid\Ga)=m_0((\theta_w)_{w\in V})$ and put
$M_v=m_0((\theta_{vw})_{w\in V})$.

\begin{lemma}\label{lem:M}
Almost surely, $M_v=\esssup(Y_v\mid\Ga)$ for every $v$, and
$M_v=\sup_i\xi_{vi}M_{vi}$. The family $(M_v)$ is an endogenous invariant recursive tree
process of \eqref{eq:hom} on $[0,\infty]$, and $M_\varnothing\ge Y_\varnothing$.
\end{lemma}

\begin{proof}
By Lemma~\ref{lem:blocks}, $\sigma(Y_v)\vee\Ga_v\subset\cF_v$ is independent of the
$\sigma$-field generated by the innovations outside the subtree of $v$, and the law of
$(Y_v,(\theta_{vw})_w)$ is that of $(Y_\varnothing,(\theta_w)_w)$. By (E3),
$\esssup(Y_v\mid\Ga)=\esssup(Y_v\mid\Ga_v)=M_v$. Since
$Y_v=\sup_i\xi_{vi}Y_{vi}$ almost surely and each $\xi_{vi}$ is $\Ga$-measurable, (E1)
and (E2) give $M_v=\sup_i\xi_{vi}M_{vi}$. The variables $(M_v:|v|=n)$ are functions of
the innovations of disjoint subtrees of depth-$n$ vertices, so they are independent,
identically distributed and independent of $(\theta_w:|w|<n)$. Thus $(M_v)$ is an invariant
recursive tree process of \eqref{eq:hom} on $[0,\infty]$, and it is $\Ga$-measurable.
Finally $Y_\varnothing\le M_\varnothing$ by definition.
\end{proof}

For $a\ge0$, $n\ge 0$ and $v\in V$ let
\[
Z_{a,v}(n)=\sup_{v=w_0,w_1,\dots,w_n}\Bigl(aY_{w_n}L_v(w_n)+\sum_{d=0}^{n-1}\eta_{w_d}L_v(w_d)\Bigr),
\]
the supremum being over paths of length $n$ starting at $v$, each $w_{j}$ a child of
$w_{j-1}$. For $v=\varnothing$ this is the variable $Z_a(n)$ of the proof of
\cite[Corollary~33]{AB}, whose law is the $n$-th iterate of the map of \eqref{eq:rde} at
$\dist(aY)$. Let $X_v$ be the variable \eqref{eq:X} computed in the subtree of $v$.

\begin{lemma}\label{lem:coupling}
Almost surely, for every $v$ the limit $Z_{a,v}=\lim_n Z_{a,v}(n)$ exists, and
\begin{equation}\label{eq:sandwich}
aY_v\le Z_{a,v}\le X_v+aY_v<\infty .
\end{equation}
Moreover $Z_{a,v}=\eta_v+\sup_i\xi_{vi}Z_{a,vi}$, and $(Z_{a,v})$ is an invariant recursive
tree process of \eqref{eq:rde} with marginal $\mu_a$.
\end{lemma}

\begin{proof}
Iterating $Y_w=\sup_i\xi_{wi}Y_{wi}$ gives $aY_v=\sup_{|w|=|v|+n,\,w\in vV}aY_wL_v(w)$, and
since $\eta\ge0$ this is at most $Z_{a,v}(n)$. Bounding the supremum of a sum by the sum of
the suprema gives $Z_{a,v}(n)\le aY_v+X_v$, because each partial sum along a path is at
most the full sum along any ray extending it. Aldous and Bandyopadhyay show that
$Z_{a,v}(n)$ is nondecreasing in $n$: given $w_n$, a child $w_{n+1}$ with
$Y_{w_n}=\xi_{w_{n+1}}Y_{w_{n+1}}$ up to an arbitrarily small error extends the path without
decreasing the bracket. This gives the limit and \eqref{eq:sandwich}. Splitting off the
first step of each path,
\[
Z_{a,v}(n)=\eta_v+\sup_i\xi_{vi}Z_{a,vi}(n-1),
\]
and since the terms increase in $n$, the suprema and the limit commute, which gives the
recursion. Each $Z_{a,v}$ is the same measurable function of
$(\theta_{vw},Y_{vw})_{w\in V}$, so by Lemma~\ref{lem:blocks} the variables
$(Z_{a,v}:|v|=n)$ are independent with a common law, independent of $(\theta_w:|w|<n)$.
The common law is $\lim_n\dist(Z_a(n))=\mu_a$.
\end{proof}

\begin{proof}[Proof of Theorem~\ref{thm:main}]
Part (a) is Lemma~\ref{lem:M}.

(b) If $(Y_v)$ is endogenous then each $Y_v$ is $\Ga$-measurable, by Lemma~\ref{lem:blocks}
applied to the shifted process. Then $Z_{a,\varnothing}(n)$, a supremum of countably many
$\Ga$-measurable variables, is $\Ga$-measurable, and so is $Z_{a,\varnothing}$. By
Lemma~\ref{lem:coupling} and uniqueness in law, $\mu_a$ is endogenous.

(c) Let $a>0$ and suppose $\mu_a$ is endogenous. The process $(Z_{a,v})$ of
Lemma~\ref{lem:coupling} has marginal $\mu_a$, so $Z_{a,\varnothing}=f(\Ga)$ almost surely
for a measurable $f$, and $f(\Ga)<\infty$ by \eqref{eq:sandwich}. Then
$Y_\varnothing\le f(\Ga)/a$ almost surely, so $M_\varnothing\le f(\Ga)/a<\infty$ almost
surely, contradicting $\Pb(M=\infty)>0$.

For the equivalence, suppose (ii) fails and let $\dist(Y^*)\ne\delta_0$ be an endogenous
solution of \eqref{eq:hom} on $[0,\infty)$. Taking $Y=Y^*$ in (b), $\mu_a$ is endogenous for
every $a$, so (i) fails. Conversely, suppose (ii) holds and let $\dist(Y)$ be any nonzero
solution. By (a), $M$ is an endogenous solution on $[0,\infty]$ with $M\ge Y_\varnothing$,
so $\dist(M)\ne\delta_0$. If $M$ were finite almost surely it would be an endogenous
solution on $[0,\infty)$ other than $\delta_0$; hence $\Pb(M=\infty)>0$, and (c) gives (i).
\end{proof}

\section{The regular case}\label{sec:regular}

We use the setting of \cite{ABM} with weights $T_i=\xi_i$. In their notation, the
assumptions are
\begin{enumerate}
\item[(A1)] $\Pb((\xi_i)\in\{0,1\}^{\N})<1$;
\item[(A2)] $\E N>1$;
\item[(A3)] there is $\alpha>0$ with $1=m(\alpha)<m(\beta)$ for all $\beta\in[0,\alpha)$;
\item[(A4a)] $m'(\alpha)=\E\sum_i\xi_i^{\alpha}\log\xi_i\in(-\infty,0)$ and
$\E\,W_1\log^+W_1<\infty$;
\item[(A4b)] $m(\theta)<\infty$ for some $\theta\in[0,\alpha)$;
\item[(A5)] there is a finite, nonnegative random variable $W$ with $\Pb(W\ne0)>0$ and
$W\stackrel{d}{=}\sum_i\xi_i^{\alpha}W_i$.
\end{enumerate}

\begin{lemma}\label{lem:assumptions}
Under the hypotheses of Theorem~\ref{thm:regular}, the point $\alpha$ is well defined and
(A1)--(A3), (A4b) and (A5) hold.
\end{lemma}

\begin{proof}
The function $m$ is convex from $[0,\infty)$ to $[0,\infty]$, finite on $[\theta_0,p]$, with
$m(\theta_0)>1>m(p)$. Hence there is $\alpha\in(\theta_0,p)$ with $m(\alpha)=1$. If
$\beta\in(\theta_0,\alpha)$ had $m(\beta)\le1$, convexity on $\beta<\alpha<p$ would give
$m(\alpha)<1$; if $\beta\in[0,\theta_0)$ had $m(\beta)\le 1$, convexity on
$\beta<\theta_0<\alpha$ would give $m(\theta_0)\le1$. So $m>1$ on $[0,\alpha)$, which is
(A3) and shows that $\alpha$ is unique; in particular $\E N=m(0)>1$, which is (A2), and
(A4b) holds with $\theta=\theta_0$. If all weights lay in $\{0,1\}$, then $m$ would be
constant, which it is not; this is (A1). For $\theta$ in a compact subinterval of
$(\theta_0,p)$, $\xi^\theta|\log\xi|$ is bounded by a constant times
$\xi^{\theta_0}+\xi^{p}$, so by dominated convergence $m$ is differentiable at $\alpha$
with $m'(\alpha)=\E\sum_i\xi_i^\alpha\log\xi_i$, and by
convexity $(m(\alpha)-m(\theta_0))/(\alpha-\theta_0)\le m'(\alpha)\le
(m(p)-m(\alpha))/(p-\alpha)<0$. So the alternative $\E W_1\log^+W_1<\infty$ of the
hypotheses gives (A4a). Now (A5) follows from \cite[Proposition~2.1]{ABM}, which requires
(A1)--(A3) together with either (A4a) or $N<\infty$ almost surely.
\end{proof}

\begin{lemma}\label{lem:unbounded}
If $m(p)<1$ for some $p>0$, every nonzero solution of \eqref{eq:hom} has unbounded
support; in particular it is not concentrated on a single point.
\end{lemma}

\begin{proof}
If $Y\le c$ almost surely, then $Y_\varnothing=\sup_{|v|=n}L(v)Y_v\le c\sup_{|v|=n}L(v)$,
and $\E\sup_{|v|=n}L(v)^p\le\E\sum_{|v|=n}L(v)^p=m(p)^n\to0$. Hence $Y_\varnothing=0$
almost surely.
\end{proof}

\begin{proof}[Proof of Theorem~\ref{thm:regular}]
Let $\dist(Y)$ be a nonzero solution of \eqref{eq:hom} with process $(Y_v)$, and put
$Y'_v=1/Y_v\in(0,\infty]$. Then $Y'_v=\inf\{Y'_{vi}/\xi_{vi}:\xi_{vi}>0\}$ with
$\inf\varnothing=\infty$, which is equation (1.3) of \cite{ABM}. Let
$f(t)=\Pb(Y'_\varnothing\ge t)$ for $t\ge0$. Then $f$ is nonincreasing with
$f(0)=f(0+)=1$, and $f(t)=\E\prod_i f(t\xi_i)$ with $f(0)=1$. If $f$ took only the values
$0$ and $1$ on $(0,\infty)$, then $Y'_\varnothing$ would be almost surely constant in
$(0,\infty]$, so $Y$ would be almost surely $0$ or almost surely a positive constant,
which is excluded by the hypothesis and by Lemma~\ref{lem:unbounded}. So $f$ belongs to the
class $\mathcal{S}(\mathcal{M})$ of \cite{ABM}.

Let $\cA_n$ be the $\sigma$-field generated by $(\xi_{vi}:|v|<n,\ i\ge1)$ and
$\cA_\infty=\bigvee_n\cA_n$. Iterating the equation,
$Y'_\varnothing=\inf\{Y'_v/L(v):|v|=n,\ L(v)>0\}$, and the variables $(Y'_v:|v|=n)$ are
independent with survival function $f$ and independent of $\cA_n$. Hence
\[
\Pb(Y'_\varnothing\ge t\mid\cA_n)=\prod_{|v|=n}f(tL(v))=:M_n(t),
\]
the multiplicative martingale of \cite[Section~8]{ABM}. By L\'evy's upward theorem,
$M_n(t)\to\Pb(Y'_\varnothing\ge t\mid\cA_\infty)$ almost surely. By Lemma~\ref{lem:assumptions}
and \cite[Theorem~8.3]{ABM}, there are an endogenous fixed point $W$ as in (A5) and a
function $h$ in the class $\mathcal{H}_r$ of \cite{ABM} (multiplicatively $r$-periodic, with
$h(t)t^\alpha$ nondecreasing) such that, for every $t\ge 0$,
\[
\Pb(Y'_\varnothing\ge t\mid\cA_\infty)=e^{-W h(t)t^{\alpha}}\quad\text{almost surely.}
\]
The function $h$ is positive: if $h(t_0)=0$, then periodicity and monotonicity of
$h(t)t^\alpha$ force $h\equiv0$, and taking expectations would give $f\equiv1$. By
Lemma~\ref{lem:eta}, conditioning on $\Ga$ is the same as conditioning on $\cA_\infty$, so
for every rational $y>0$,
\[
\Pb(Y_\varnothing>y\mid\Ga)=1-e^{-W h(1/y)y^{-\alpha}}>0\qquad\text{almost surely on }\{W>0\}.
\]
Hence $M=\infty$ almost surely on $\{W>0\}$, an event of positive probability, and
Theorem~\ref{thm:main}(c) completes the proof.
\end{proof}

\section{Examples and remarks}\label{sec:remarks}

\begin{example}
Let every vertex have four children with weight $1/2$. Then $m(\theta)=4\cdot2^{-\theta}$
and $\alpha=2$. The Fr\'echet law $\Pb(Y\le y)=e^{-y^{-2}}$ solves \eqref{eq:hom}, and
$(Y_v)$ is independent of $\Ga$, so $M=\infty$ and, for any bounded $\eta\ge0$, no $\mu_a$
with $a>0$ is endogenous. In contrast, with two children of weight $1$ and $\eta=0$ the
constant $Y=1$ is an endogenous solution and $\mu_a=\delta_a$ is endogenous for every $a$,
as Theorem~\ref{thm:main}(b) predicts; here $m\equiv2$ and Theorem~\ref{thm:regular} does
not apply.
\end{example}

\begin{example}
In Example~37 of \cite{AB}, $\xi_i=c\in(0,1)$ for $i\le N$, where $N$ is the offspring
variable of a Galton--Watson tree. If $1<\E N<\infty$ then $m(\theta)=c^\theta\E N$, so
Theorem~\ref{thm:regular} applies with $\theta_0=0$ and $\alpha=\log\E N/\log(1/c)$. In
Example~38 of \cite{AB}, $\xi_i=e^{-\zeta_i}$ with $(\zeta_i)$ the points of a rate one
Poisson process; then $N=\infty$, $m(\theta)=1/\theta$, $\alpha=1$, and $W_1=\sum_ie^{-\zeta_i}$
has finite second moment, so the hypotheses hold. In both examples every invariant measure
$\mu_a$, $a>0$, built from a nonzero solution of \eqref{eq:hom} is non-endogenous.
\end{example}

\begin{remark}
In the regular case the dichotomy in Theorem~\ref{thm:main} does not leave room for an
intermediate situation: the process of $Y$ is never endogenous and $M$ is infinite with
positive probability. Outside the regular case, a counterexample to Conjecture~\ref{conj}
must come from an endogenous solution of \eqref{eq:hom} on $[0,\infty)$ other than
$\delta_0$, and when $m(p)<1$ for some $p>0$ such a solution must have unbounded support by
Lemma~\ref{lem:unbounded}. Two cases remain open: weights without a characteristic
exponent, where $m$ jumps from $\infty$ to a value at most $1$ (the case studied for the
min-type equation by Jagers and R\"osler \cite{JR} and Alsmeyer and Meiners \cite{AM}), and
infinitely many positive weights for which $\E W_1\log^+W_1=\infty$.
\end{remark}

\begin{remark}
The correspondence between solutions of an inhomogeneous max-type equation and its
homogeneous counterpart was studied by R\"uschendorf \cite{R}. Theorem~\ref{thm:main}
transfers endogeny along the correspondence of \cite[Corollary~33]{AB} in one direction
and transfers its failure through the conditional supremum $M$ in the other.
\end{remark}

\section*{Acknowledgments and funding}

This research was supported by the European Research Council (ERC) under the European Union's Horizon Europe research and innovation programme (grant agreement No.~101041711), by the Simons Foundation as part of the Collaboration on the Mathematical and Scientific Foundations of Deep Learning, by Heights Labs, by Convex Nexus Capital, by the Israel Science Foundation (grant number 2258/19), and by the Israel Science Foundation (ISF Grant 4101/25).

\section*{Declaration of competing interest}

The author holds equity in Heights Labs and in Convex Nexus Capital and has a management role at Convex Nexus Capital. These entities are included in the funding declaration above.

\section*{Declaration of generative AI and AI-assisted technologies}

Anthropic Claude AI systems were used to develop the mathematical argument, check the cited results, and write and prepare this manuscript and its publication files. The author selected the problem and directed the work; the author did not carry out the mathematics or the writing. Responsibility for publication rests with the author. The AI-assisted internal reviews and compilation checks are not external peer review or formal verification.

\section*{Data and code availability}

No data or code were used. The source of this manuscript is in the repository
\[
\texttt{github.com/yspennstate/discounted-tree-sums-endogeny}.
\]

\begin{thebibliography}{9}

\bibitem{AB}
D.~J.~Aldous and A.~Bandyopadhyay, A survey of max-type recursive distributional
equations, \emph{Ann. Appl. Probab.} 15 (2005), 1047--1110.

\bibitem{ABM}
G.~Alsmeyer, J.~D.~Biggins and M.~Meiners, The functional equation of the smoothing
transform, \emph{Ann. Probab.} 40 (2012), 2069--2105.

\bibitem{AM}
G.~Alsmeyer and M.~Meiners, A min-type stochastic fixed-point equation related to the
smoothing transformation, \emph{Theory Stoch. Process.} 15 (2009), 19--41.

\bibitem{JR}
P.~Jagers and U.~R\"osler, Stochastic fixed points for the maximum, in
\emph{Mathematics and Computer Science III}, Trends Math., Birkh\"auser, Basel, 2004,
325--338.

\bibitem{R}
L.~R\"uschendorf, On stochastic recursive equations of sum and max type, \emph{J. Appl.
Probab.} 43 (2006), 687--703.

\end{thebibliography}

\end{document}
